\section{An involution and a sign}\label{sec:involution}

\subsection{The finite set and its maps}
Let $p$ be as in \cref{thm:cone-parity}, and let $D$ be a positive nonsquare with
$p\mid D$. First assume $D\equiv1\pmod{40}$, and put
\[
  \Tset_p(D)=\{(x,y,z)\in\Z_{>0}^3:\ x^2+2yz=D,\ \ p\nmid\gcd(x,y,z)\}.
\]
The last condition says only that $p$ does not divide all three coordinates.
The set is finite, its $x$-coordinates are odd, and $4\mid yz$.

Each triple $t=(x,y,z)$ falls into one of five \emph{branches}
$\beta(t)\in\{0,1,2,3,4\}$, according to the position of $2y$ relative to the
increasing numbers $x<2x+z<3x+2z<4x+4z$ (\cref{tab:jackson}). On branch $b$
\emph{Jackson's map} $\Jk$ is the linear map in the last column; see \cite{Jackson}
for the classical version, which is stated for prime discriminants.

\begin{table}[ht]
\caption{Branches of Jackson's map $\Jk$.}\label{tab:jackson}
\begin{tabular}{cll}
\toprule
$\beta$ & condition & $\Jk(x,y,z)$\\
\midrule
0 & $2y<x$ & $(x-2y,\ z+2x-2y,\ y)$\\
1 & $x<2y<2x+z$ & $(2y-x,\ y,\ 2x-2y+z)$\\
2 & $2x+z<2y<3x+2z$ & $(3x-2y+2z,\ 2x-y+2z,\ -2x+2y-z)$\\
3 & $3x+2z<2y<4x+4z$ & $(-3x+2y-2z,\ -2x+2y-z,\ 2x-y+2z)$\\
4 & $4x+4z<2y$ & $(x+2z,\ z,\ -2x+y-2z)$\\
\bottomrule
\end{tabular}
\end{table}

No triple lies on a boundary: $2y=x$ and $2y=3x+2z$ contradict the oddness of
$x$, while $2y=2x+z$ and $2y=4x+4z$ would give $D=(x+z)^2$ and $D=(x+2z)^2$. A
direct substitution shows that $\Jk$ preserves $x^2+2yz$, produces positive
coordinates, and sends branch $b$ to branch $(4,1,2,3,0)_b$. The linear maps on
branches $0$ and $4$ are inverse to each other, and those on branches $1,2,3$
are involutions; all have integral inverses, so $p\nmid\gcd$ is preserved. Hence
$\Jk$ is an involution of $\Tset_p(D)$.

The swap $\Sw(x,y,z)=(x,z,y)$ is a second involution. For $\ell\in\{2,5\}$
define the rational map
\[
  A_\ell(x,y,z)=(x,\ y/\ell,\ \ell z).
\]
It restricts to a bijection from the triples of $\Tset_p(D)$ with $\ell\mid y$ to
those with $\ell\mid z$; the condition $p\nmid\gcd$ is preserved because
$p\nmid10$.

\subsection{The sign}
Write $\chi(a)=\Leg ap$. By \eqref{eq:legendre-values},
$\chi(-1)=1$, $\chi(2)=-1$, $\chi(5)=1$. Let $\F_2=\Z/2\Z$. For
$t=(x,y,z)\in\Tset_p(D)$ define a sign $\epsilon(t)$ and a bit $g(t)\in\F_2$ by
\[
  \epsilon(x,y,z)=\begin{cases}\chi(y),&p\nmid y,\\ \chi(-2z),&p\mid y,\end{cases}
  \qquad (-1)^{g(t)}=\epsilon(t).
\]
This is well defined: if $p\mid y$ then $p\mid x$ (as $p\mid D$), so the
condition $p\nmid\gcd$ forces $p\nmid z$.

\begin{lemma}\label{lem:sign}
In $\F_2$, on the domains of the respective maps,
\[
  g(\Sw t)=g(t)+1,\qquad g(\Jk t)=g(t)+c_{\beta(t)},\quad (c_0,\dots,c_4)=(1,0,0,1,1),
\]
\[
  g(A_2t)=g(t)+1,\qquad g(A_5t)=g(t).
\]
\end{lemma}

\begin{proof}
Suppose first that $p\nmid y$. Since $p\mid D$, $z\equiv-x^2/(2y)\pmod p$. In
branches $0,2,3$ the new second coordinate is then congruent to
\[
  -\frac{(x-2y)^2}{2y},\qquad -\frac{(x-y)^2}{y},\qquad \frac{(x-2y)^2}{2y}
\]
respectively; in branch $1$ it is $y$, and in branch $4$ it is $z\equiv-x^2/(2y)$.
Since $\chi(-1)=1$ and $\chi(2)=-1$, these give the stated changes whenever the new
second coordinate is a unit modulo $p$. If it is divisible by $p$, use the new third
coordinate in the definition of $\epsilon$: in branches $0,2,3,4$ this gives
$\chi(-2y)$, $\chi(-y)$, $\chi(2y)$, $\chi(-2y)$ respectively (at the roots
$x\equiv2y$, $y$, $2y$, $0$), the same multipliers of $\chi(y)$.

If $p\mid y$, then $p\mid x$ and $p\nmid z$. The new second coordinates in
branches $0,2,3,4$ are congruent to $z,2z,-z,z$, and in branch $1$ the new third
coordinate is $z$; comparing with $\chi(-2z)$ gives the same table.

For $\Sw$: if $y,z$ are both units then $\chi(yz)=\chi(-x^2/2)=\chi(-2)=-1$, so
$\chi(z)=-\chi(y)$; if one of them is divisible by $p$, the second case of the
definition gives the same reversal. Finally $A_\ell$ multiplies the relevant
coordinate by $\chi(\ell)^{\pm1}$, giving $+1$ for $\ell=2$ and $0$ for $\ell=5$.
\end{proof}
